\documentclass[aspectratio=169]{beamer}

\usepackage[T1]{fontenc}
\usepackage[utf8]{inputenc}
\usepackage[english,brazil]{babel}
\usepackage{lmodern}
\usepackage{amsmath,amssymb,booktabs,array}
\usepackage{tikz}
\usepackage{pgfplots}
\usepackage{siunitx}
\usepackage{ragged2e}
\usepackage{hyperref}

\usetikzlibrary{positioning,arrows.meta,calc,fit,backgrounds}
\pgfplotsset{compat=1.18}
\sisetup{locale = DE}

\usetheme{ifsclean}

\definecolor{ifscgreen}{RGB}{0,128,0}

\title{Análise do Showcase Queueing Time}
\subtitle{Avaliação de Desempenho com OMNeT++/INET e apoio de IA}
\author{Gabriel Kuster Kraus e Wagner Flores dos Santos}
\institute{ADS29009 -- Avaliação de Desempenho de Sistemas \\ Engenharia de Telecomunicações}
\date{Junho de 2026}

\begin{document}

% ---------------------------------------------------------
\begin{frame}[plain,noframenumbering]
    \maketitle

    \vspace{0.5cm}
    \centering
    \small
    Showcase de referência: \\
    \url{https://inet.omnetpp.org/docs/showcases/measurement/queueingtime/doc/index.html}
\end{frame}

% ---------------------------------------------------------
\begin{frame}{Agenda}
  \tableofcontents
\end{frame}

% =========================================================
\section{Introdução}

\begin{frame}{Contexto e objetivo}
\begin{itemize}
    \item O trabalho analisa um experimento já preparado de avaliação de desempenho de redes usando OMNeT++ e INET.
    \item O showcase escolhido é o \textbf{Queueing Time}, voltado à medição do tempo de espera de pacotes em fila.
    \item O objetivo principal é compreender como a carga da rede influencia o atraso, a formação de filas e o desempenho geral do sistema.
    \item Além da análise do showcase, é proposta uma variação experimental com coleta de dados e apoio de IA.
\end{itemize}

\begin{informacao}
O foco da apresentação é relacionar o modelo de simulação com métricas clássicas de Avaliação de Desempenho de Sistemas.
\end{informacao}
\end{frame}

% ---------------------------------------------------------
\begin{frame}{Breve revisão de conceitos}
\begin{itemize}
    \item \textbf{Fila:} estrutura onde pacotes aguardam atendimento quando o servidor está ocupado.
    \item \textbf{Tempo de fila:} tempo que um pacote permanece esperando antes de ser processado.
    \item \textbf{Taxa de chegada:} ritmo em que pacotes entram no sistema.
    \item \textbf{Taxa de serviço:} ritmo em que pacotes são processados.
    \item \textbf{Saturação:} ocorre quando a taxa de chegada se aproxima ou supera a capacidade de atendimento.
\end{itemize}

\begin{informacao}
Quando a carga aumenta, o tempo de espera tende a crescer de forma significativa.
\end{informacao}
\end{frame}

% =========================================================
\section{Modelo simulado}

\begin{frame}{Modelo básico do sistema}
\centering
\begin{tikzpicture}[
    >=Latex,
    bloco/.style={
        draw,
        rounded corners,
        minimum width=2.7cm,
        minimum height=1.0cm,
        align=center,
        font=\small
    }
]

\node[bloco, fill=green!10] (source) {Source\\Geração de pacotes};
\node[bloco, fill=yellow!15, right=1.0cm of source] (queue) {Queue\\Fila};
\node[bloco, fill=blue!10, right=1.0cm of queue] (server) {Server\\Atendimento};
\node[bloco, fill=gray!15, right=1.0cm of server] (sink) {Sink\\Destino};

\draw[->, thick] (source) -- (queue);
\draw[->, thick] (queue) -- (server);
\draw[->, thick] (server) -- (sink);

\end{tikzpicture}

\vspace{0.8cm}

\begin{itemize}
    \item O sistema representa uma sequência simples de geração, enfileiramento, atendimento e recepção dos pacotes.
    \item A fila é o principal ponto de interesse, pois nela ocorre o tempo de espera analisado.
\end{itemize}
\end{frame}

% ---------------------------------------------------------
\begin{frame}{Estrutura dos módulos no INET}
\begin{columns}[T,onlytextwidth]

\column{0.48\textwidth}
\textbf{Módulos principais}
\begin{itemize}
    \item \texttt{UdpSourceApp}
    \item \texttt{EthernetInterface}
    \item \texttt{PacketQueue}
    \item \texttt{UdpSinkApp}
\end{itemize}

\column{0.52\textwidth}
\textbf{Função de cada módulo}
\begin{itemize}
    \item Source: gera pacotes UDP.
    \item Queue: armazena pacotes temporariamente.
    \item Interface Ethernet: transmite os pacotes.
    \item Sink: recebe os pacotes no destino.
\end{itemize}

\end{columns}

\vspace{0.5cm}
\begin{informacao}
Esse modelo permite observar diretamente o impacto da taxa de chegada sobre o tempo de fila e a vazão.
\end{informacao}
\end{frame}

% =========================================================
\section{Configuração da simulação}

\begin{frame}{Configuração da rede}
\begin{table}
\centering
\small
\begin{tabular}{ll}
\toprule
\textbf{Elemento} & \textbf{Descrição} \\
\midrule
Ambiente & OMNeT++ 6.4.0 com INET 4.6.0 \\
Modelo & QueueingTimeMeasurementShowcase \\
Tempo de simulação & \SI{5}{s} \\
Aplicação de origem & \texttt{UdpSourceApp} \\
Aplicação de destino & \texttt{UdpSinkApp} \\
Tamanho do pacote & \SI{1200}{B} \\
Métrica principal & Tempo de fila dos pacotes \\
Saída & Arquivos \texttt{.sca}, \texttt{.vec} e CSV \\
\bottomrule
\end{tabular}
\end{table}

\vspace{0.4cm}

\begin{itemize}
    \item A configuração original utiliza geração de pacotes com intervalo exponencial.
    \item O comentário do próprio \texttt{omnetpp.ini} indica tráfego de origem próximo de \SI{96}{Mbps}.
\end{itemize}
\end{frame}

% ---------------------------------------------------------
\begin{frame}{Fatores, níveis e parâmetros fixados}
\begin{table}
\centering
\small
\begin{tabular}{lll}
\toprule
\textbf{Tipo} & \textbf{Parâmetro} & \textbf{Valores} \\
\midrule
Fator & Intervalo de geração & 100, 50, 25 e 10 us \\
Parâmetro fixado & Tempo de simulação & \SI{5}{s} \\
Parâmetro fixado & Tamanho do pacote & \SI{1200}{B} \\
Parâmetro fixado & Modelo da rede & Showcase original \\
Métrica & Throughput & Mbps \\
Métrica & Queueing Time & us / ms \\
\bottomrule
\end{tabular}
\end{table}

\vspace{0.5cm}

\begin{informacao}
Na variação proposta, o principal fator analisado foi o intervalo de geração dos pacotes.
\end{informacao}
\end{frame}

% =========================================================
\section{Coleta de estatísticas}

\begin{frame}{Métricas analisadas}
\begin{itemize}
    \item \textbf{Queueing Time:} tempo que cada pacote permanece aguardando na fila.
    \item \textbf{Queue Length:} quantidade de pacotes presentes na fila ao longo do tempo.
    \item \textbf{Throughput:} taxa efetiva de transmissão observada.
    \item \textbf{Packet Loss:} perdas ocasionadas por limitação de fila ou sobrecarga.
\end{itemize}

\begin{informacao}
Essas métricas permitem avaliar se o sistema está operando em condição estável ou próximo da saturação.
\end{informacao}
\end{frame}

% ---------------------------------------------------------
\begin{frame}[fragile]{Exportação dos resultados}
Os resultados do OMNeT++ foram exportados para análise externa usando o \texttt{opp\_scavetool}.

\vspace{0.4cm}
\footnotesize
\begin{block}{Comando utilizado}
\begin{verbatim}
opp_scavetool export -F CSV-R -o cenarios.csv results/*.sca results/*.vec
\end{verbatim}
\end{block}

\vspace{0.4cm}

\begin{itemize}
    \item Os arquivos \texttt{.sca} armazenam estatísticas escalares.
    \item Os arquivos \texttt{.vec} armazenam séries temporais.
    \item O CSV consolidado foi analisado com Python, Pandas e Matplotlib.
\end{itemize}
\end{frame}

% =========================================================
\section{Resultados do showcase}

\begin{frame}{Execução da simulação original}
\begin{itemize}
    \item O showcase original foi executado em linha de comando com \texttt{Cmdenv}.
    \item A simulação processou aproximadamente \textbf{695.444 eventos}.
    \item Foram criadas aproximadamente \textbf{198.100 mensagens/pacotes}.
    \item O tempo simulado foi de \SI{5}{s}.
    \item O throughput observado ficou próximo de \SI{99}{Mbps}.
\end{itemize}

\begin{informacao}
Esses valores confirmam que o experimento foi executado e que os dados foram coletados para análise externa.
\end{informacao}
\end{frame}

% ---------------------------------------------------------
\begin{frame}{Análise esperada do Queueing Time}
\begin{itemize}
    \item Em baixa carga, os pacotes encontram pouca ou nenhuma fila.
    \item Com carga intermediária, o tempo de fila começa a crescer.
    \item Quando a taxa de chegada se aproxima da taxa de serviço, a fila tende a aumentar rapidamente.
    \item Em condição de sobrecarga, o sistema pode apresentar atraso elevado e perdas.
\end{itemize}

\begin{informacao}
O tempo de fila é uma métrica sensível à relação entre chegada e atendimento.
\end{informacao}
\end{frame}

% ---------------------------------------------------------
\begin{frame}{Crítica sobre a apresentação dos resultados}
\begin{itemize}
    \item O showcase apresenta bem o funcionamento básico da métrica de tempo de fila.
    \item Porém, a análise poderia ser ampliada com mais repetições e sementes diferentes.
    \item Também seria interessante apresentar intervalo de confiança.
    \item A comparação entre diferentes níveis de carga poderia ser destacada de forma mais direta.
    \item A análise externa com Pandas facilita a geração de gráficos comparativos.
\end{itemize}
\end{frame}

% =========================================================
\section{Variação proposta}

\begin{frame}{Objetivo da variação experimental}
\begin{itemize}
    \item Propor uma nova perspectiva de análise sobre o showcase original.
    \item Avaliar o impacto do aumento progressivo da taxa de chegada de pacotes.
    \item Identificar o ponto em que o sistema começa a apresentar crescimento acentuado do tempo de fila.
    \item Comparar tempo de fila e vazão em diferentes cenários.
\end{itemize}

\begin{informacao}
A variação transforma o showcase em um pequeno estudo de sensibilidade da carga sobre o desempenho.
\end{informacao}
\end{frame}

% ---------------------------------------------------------
\begin{frame}{Cenários simulados}
\begin{table}
\centering
\small
\begin{tabular}{cccc}
\toprule
\textbf{Cenário} & \textbf{Intervalo de geração} & \textbf{Carga relativa} & \textbf{Expectativa} \\
\midrule
C1 & 100 us & Original & Fila moderada \\
C2 & 50 us & 2x maior & Maior atraso \\
C3 & 25 us & 4x maior & Saturação \\
C4 & 10 us & 10x maior & Saturação acentuada \\
\bottomrule
\end{tabular}
\end{table}

\vspace{0.4cm}

\begin{informacao}
A comparação entre cenários permite observar a transição entre operação controlada e saturação.
\end{informacao}
\end{frame}

% =========================================================
\section{Resultados obtidos}

\begin{frame}{Resultados obtidos -- Throughput}
\centering
\includegraphics[width=0.70\linewidth]{figs/grafico_throughput_cenarios.png}

\footnotesize
\begin{itemize}
    \item O throughput permaneceu próximo de \SI{99}{Mbps} em todos os cenários.
    \item O aumento da carga de entrada não gerou aumento proporcional da vazão.
    \item Isso indica que o enlace passou a operar próximo de sua capacidade máxima.
\end{itemize}
\end{frame}

% ---------------------------------------------------------
\begin{frame}{Resultados obtidos -- Tempo em fila}
\centering
\includegraphics[width=0.70\linewidth]{figs/grafico_queueingtime_cenarios.png}

\footnotesize
\begin{itemize}
    \item O cenário C1 apresentou tempo médio em fila de aproximadamente \SI{51}{ms}.
    \item A partir de C2, o tempo médio ficou próximo de \SI{100}{ms}.
    \item O aumento da carga não elevou a vazão, mas aumentou o atraso.
    \item Esse comportamento caracteriza saturação do sistema.
\end{itemize}
\end{frame}

% ---------------------------------------------------------
\begin{frame}{Resumo quantitativo dos cenários}
\begin{table}
\centering
\small
\begin{tabular}{cccc}
\toprule
\textbf{Cenário} & \textbf{Intervalo} & \textbf{Throughput médio} & \textbf{Tempo médio em fila} \\
\midrule
C1 & 100 us & \SI{98,90}{Mbps} & \SI{51,08}{ms} \\
C2 & 50 us  & \SI{99,05}{Mbps} & \SI{99,18}{ms} \\
C3 & 25 us  & \SI{99,05}{Mbps} & \SI{99,89}{ms} \\
C4 & 10 us  & \SI{99,05}{Mbps} & \SI{100,13}{ms} \\
\bottomrule
\end{tabular}
\end{table}

\vspace{0.4cm}

\begin{informacao}
O throughput estabiliza, enquanto o tempo em fila cresce. Esse é o principal indicativo de saturação.
\end{informacao}
\end{frame}

% ---------------------------------------------------------
\begin{frame}{Análise dos resultados}
\begin{itemize}
    \item O aumento da taxa de geração elevou a pressão sobre a fila.
    \item A vazão não aumentou porque o sistema já estava próximo do limite de transmissão.
    \item O impacto do aumento de carga apareceu principalmente no \textbf{tempo em fila}.
    \item A diferença mais significativa ocorreu entre C1 e C2.
    \item A partir de C2, os cenários permaneceram próximos de \SI{100}{ms}, indicando regime de saturação.
\end{itemize}

\begin{informacao}
Em sistemas saturados, aumentar a chegada de pacotes tende a aumentar atraso e perdas, não necessariamente a vazão.
\end{informacao}
\end{frame}
\begin{frame}[fragile]{Dataset Gerado}

\small

\begin{itemize}
\item Arquivos produzidos pelo OMNeT++:
\begin{itemize}
\item General-\#0.sca
\item General-\#0.vec
\end{itemize}

\item Conversão para CSV usando:

\texttt{opp\_scavetool}

\item Dataset final:

\textbf{cenarios.csv (241 MB)}
\end{itemize}

\vspace{0.3cm}

\begin{verbatim}
Exported 172 scalars,
1238 parameters,
105 vectors,
20 histograms
\end{verbatim}

\end{frame}

\begin{frame}[fragile]{Exemplo de Dados Exportados}

\scriptsize

\begin{verbatim}
run,type,module,name,mean,max

C1,histogram,...,
queueingTime:histogram,
0.051075,0.093588

C2,histogram,...,
queueingTime:histogram,
0.099177,0.101279

C3,histogram,...,
queueingTime:histogram,
0.099887,0.101279

C4,histogram,...,
queueingTime:histogram,
0.100129,0.101279
\end{verbatim}

\vspace{0.2cm}

Os dados acima foram extraídos diretamente
do arquivo CSV exportado pelo OMNeT++.

\end{frame}

\begin{frame}[fragile]{Extração das Métricas com Python}

\scriptsize

\begin{verbatim}
df = pd.read_csv("cenarios.csv")

q = df[
 (df["type"]=="histogram") &
 (df["name"]=="queueingTime:histogram")
]

q["cenario"] = q["run"].str.extract(r"^(C\d)")
q["mean_us"] = q["mean"] * 1e6
\end{verbatim}

\vspace{0.3cm}

O script foi utilizado para extrair as
métricas do dataset e gerar os gráficos.

\end{frame}


% =========================================================
\section{Dataset e análise externa}

\begin{frame}{Fluxo de geração do dataset}
\centering
\begin{tikzpicture}[
    >=Latex,
    bloco/.style={
        draw,
        rounded corners,
        minimum width=2.6cm,
        minimum height=0.9cm,
        align=center,
        font=\small
    }
]

\node[bloco, fill=green!10] (omnet) {OMNeT++\\INET};
\node[bloco, fill=yellow!15, right=0.8cm of omnet] (sca) {.sca / .vec};
\node[bloco, fill=blue!10, right=0.8cm of sca] (csv) {CSV\\scavetool};
\node[bloco, fill=gray!15, right=0.8cm of csv] (pandas) {Pandas\\Gráficos};

\draw[->, thick] (omnet) -- (sca);
\draw[->, thick] (sca) -- (csv);
\draw[->, thick] (csv) -- (pandas);

\end{tikzpicture}

\vspace{0.8cm}

\begin{itemize}
    \item O OMNeT++ gera os arquivos de resultado.
    \item O \texttt{opp\_scavetool} converte os dados para CSV.
    \item O Python permite filtrar métricas, agrupar cenários, calcular médias e gerar gráficos.
\end{itemize}
\end{frame}

% ---------------------------------------------------------
\begin{frame}{Dataset gerado}
\begin{itemize}
    \item Foi gerado um arquivo CSV consolidado com aproximadamente \textbf{241 MB}.
    \item O arquivo contém dados escalares, vetoriais, histogramas e parâmetros de configuração.
    \item As métricas extraídas para os gráficos foram:
    \begin{itemize}
        \item \texttt{throughput:histogram}
        \item \texttt{queueingTime:histogram}
    \end{itemize}
    \item O dataset pode ser versionado no GitHub junto com os scripts de análise.
\end{itemize}

\begin{informacao}
Para reprodutibilidade, recomenda-se publicar também o \texttt{omnetpp.ini}, os scripts Python e os gráficos gerados.
\end{informacao}
\end{frame}

% ---------------------------------------------------------
\begin{frame}[fragile]{Exemplo de análise com Pandas}
\begin{block}{Trecho usado para extrair o tempo médio em fila}
\begin{verbatim}
df = pd.read_csv("cenarios.csv")

q = df[
    (df["type"] == "histogram") &
    (df["name"] == "queueingTime:histogram")
].copy()

q["cenario"] = q["run"].str.extract(r"^(C\d)")
q["mean_ms"] = q["mean"].astype(float) * 1e3
\end{verbatim}
\end{block}

\begin{itemize}
    \item O processamento externo permitiu comparar os cenários de forma direta.
\end{itemize}
\end{frame}

% =========================================================
\section{Uso de IA em ADS}

\begin{frame}{Uso da IA no projeto}
\begin{itemize}
    \item A IA foi utilizada para auxiliar na instalação e configuração do ambiente OMNeT++.
    \item Auxiliou na execução do showcase Queueing Time.
    \item Foi utilizada para criar cenários adicionais de carga.
    \item Auxiliou na extração das métricas dos arquivos \texttt{.sca} e \texttt{.vec}.
    \item Foi utilizada para gerar scripts Python para análise dos resultados.
    \item Auxiliou na interpretação dos gráficos e identificação do ponto de saturação do sistema.
\end{itemize}
\end{frame}

% ---------------------------------------------------------
\begin{frame}[fragile]{Exemplos de prompts usados}

\small
\begin{block}{Prompts possíveis}
\begin{verbatim}
Como exportar os resultados do OMNeT++ para CSV usando scavetool?

Localize no CSV as métricas relacionadas a queueing time.

Gere um script Python para comparar throughput entre os cenários.

Explique por que o throughput estabilizou, mas o tempo em fila aumentou.
\end{verbatim}
\end{block}

\vspace{0.3cm}  

\begin{itemize}
    \item Os prompts ajudaram na exploração dos resultados e na construção da análise.
\end{itemize}
\end{frame}

% ---------------------------------------------------------
\begin{frame}{Relação com o showcase Aloha AI}
\begin{itemize}
    \item O showcase Aloha AI demonstra uma interação em linguagem natural com a simulação.
    \item No contexto de ADS, essa abordagem pode ser usada para investigar eventos críticos.
    \item Em vez de procurar colisões, como no Aloha, o foco foi encontrar cenários de maior atraso ou saturação.
    \item Assim, o estudo com INET permanece como base principal, e o Aloha AI entra como inspiração para a parte de IA.
\end{itemize}

\vspace{0.4cm}

\centering
\small
\url{https://omnetpp.org/documentation/ai-demos/aloha/showcase}
\end{frame}

% =========================================================

\section{Conclusão}

\begin{frame}{Conclusões}
\begin{itemize}
    \item O showcase Queueing Time é adequado para estudar conceitos de filas e atraso em redes.
    \item A estrutura do INET permite coletar estatísticas relevantes para Avaliação de Desempenho.
    \item A variação proposta avaliou o impacto da carga sobre o tempo de fila e a vazão.
    \item Os resultados mostraram throughput praticamente constante e aumento significativo do tempo em fila.
    \item O uso de IA auxiliou na exploração dos cenários, extração das métricas e interpretação dos resultados.
\end{itemize}
\end{frame}

% ---------------------------------------------------------
\begin{frame}{Possibilidades futuras}
\begin{itemize}
    \item Executar múltiplas repetições com sementes diferentes.
    \item Calcular intervalo de confiança.
    \item Avaliar diferentes tamanhos de fila.
    \item Comparar diferentes políticas de atendimento.
    \item Usar IA para sugerir automaticamente novos cenários de simulação.
\end{itemize}

\begin{informacao}
A principal evolução seria transformar o showcase em um experimento completo com análise estatística.
\end{informacao}
\end{frame}

% =========================================================
\begin{frame}{Referências}

\footnotesize

\begin{thebibliography}{99}

\bibitem{inet}
INET Framework. Queueing Time Measurement Showcase.
Disponível em: \url{https://inet.omnetpp.org/docs/showcases/measurement/queueingtime/doc/index.html}.

\bibitem{omnet}
OMNeT++ Discrete Event Simulator.
Disponível em: \url{https://omnetpp.org}.

\bibitem{repo}
SANTOS, Wagner F.; equipe.
Projeto ADS 2026.1 – Avaliação A3.2.
Disponível em:
\url{https://github.com/wagnerfloresdossantos/engtelecom-ifsc-sj/tree/main/ads-2026.1/a3.2}

\end{thebibliography}

\end{frame}

\end{document}

